% Endpoint construction for the proof of Theorem 4.2.
% Assumes c=(sqrt(3)-1)/2 and c<=lambda<1/2.
\paragraph{Nonempty endpoint fibres.}
We first construct infinitely many nontrivial points of the intersection for
rational $\alpha>1$. The construction will also give the strict margins needed
below. Put
\[
 A:=\log\bigl((1-\lambda)+\lambda C_\lambda\bigr),
 \qquad b:=-\log(1-\lambda),\qquad h:=-\log\lambda.
\]
Thus $\operatorname{conv}A=[-b,0]$ and $0<b<h$.
We will show that, for every real $\alpha>1$, there is an integer $j$ such that
\begin{equation}\label{eq:coarse-strict-link}
 \alpha A\cap
 (\log\alpha-jh-b,\log\alpha-jh)\ne\varnothing.
\end{equation}
This is a statement about a point lying strictly inside a translated hull;
we will use the Gap Lemma to obtain actual intersections of Cantor sets.

First, $A$ has thickness greater than one. Consider a cylinder hull $[a,a+d]$
inside $[1-\lambda,1]$, where $d=\lambda^n$ and $n\ge1$.
Its central gap is $(a+\lambda d,a+(1-\lambda)d)$.
Since the derivative of the logarithm decreases, the logarithmic left bridge
is at least $\lambda/(1-2\lambda)>1$ times the logarithmic gap.
The right bridge is longer than the gap precisely when
\[
 (a+d)(a+\lambda d)>(a+(1-\lambda)d)^2,
\]
or, equivalently,
\[
 (3\lambda-1)+(3\lambda-1-\lambda^2)\frac da>0.
\]
Here $0<d/a\le\lambda/(1-\lambda)$. If the second coefficient is negative,
the left side is at least
\[
 \frac{3\lambda-1-\lambda^3}{1-\lambda}>0;
\]
the numerator increases on $[c,1/2]$ and its value at $c$ is $(3c-1)/2>0$.
If the coefficient is nonnegative, positivity is immediate.
The bridge-to-gap ratios extend continuously to $d/a=0$, with limiting
value $\lambda/(1-2\lambda)>1$, so they have a strict common lower bound
above one for each fixed $\lambda$.
Moreover, every logarithmic gap one generation deeper is at most
$\lambda/(1-\lambda)<1$ times every gap of the preceding generation:
this follows by integrating $1/x$ over the corresponding gaps in
$[1-\lambda,1]$. Hence later gaps cannot truncate these bridges when the
gaps are removed in decreasing order. This proves the thickness assertion.

The Gap Lemma~\cite{Newhouse1979} now gives the exact identity
\begin{equation}\label{eq:coarse-thickening}
 \alpha A-A=\alpha A+[0,b],\qquad \alpha\ge1.
\end{equation}
Indeed, if $s\in\alpha A+[0,b]$, the hull of $s+A$ meets $\alpha A$.
A contact at a hull endpoint already gives an intersection.
Otherwise $s+A$ is not contained in a gap of $\alpha A$, because its hull
meets $\alpha A$, and $\alpha A$ is not contained in a gap of $s+A$,
because its diameter is at least $b$, larger than every bounded gap of $A$;
the hull intersection excludes containment in an unbounded gap.
The thickness product exceeds one, so the two Cantor sets intersect.
The reverse inclusion is immediate from $-A\subset[0,b]$.
The largest gap of $A$ has length
\[
 \log\frac{1-\lambda^2}{1-\lambda+\lambda^2}.
\]
Consequently, if this length multiplied by $\alpha$ is strictly less than
$b$, then $\alpha A+[0,b]=[-\alpha b,b]$, with strict overlap across each
gap. If also $(\alpha+1)b>h$, every residue modulo $h$ has a representative
$s$ for which $\alpha A\cap(s-b,s)\ne\varnothing$.
More generally, the same conclusion holds whenever every complementary arc
of $\alpha A$ modulo $h$ has length strictly less than $b$.

For $1<\alpha\le23/19$, we have
$0<\log\alpha<b$, since $23/19<3/2<1/(1-c)$.
The point $0\in\alpha A$ therefore proves
\eqref{eq:coarse-strict-link} with $j=0$.
For $23/19\le\alpha\le3$, the ratio
$(1-\lambda^2)/(1-\lambda+\lambda^2)$ decreases with $\lambda$:
its derivative has the sign of $1-4\lambda+\lambda^2<0$.
Using $\lambda\ge c>4/11$ gives
\[
 \left(\frac{1-\lambda^2}{1-\lambda+\lambda^2}\right)^3(1-\lambda)
 <\left(\frac{35}{31}\right)^3\frac7{11}
 =\frac{300125}{327701}<1.
\]
Thus all gaps of $\alpha A$ are shorter than $b$.
Also $(42/19)b>h$. It suffices to verify this at $\lambda=c$;
using $1-c=\sqrt3c$, it is equivalent to $3^{21}c^{23}<1$.
Since $c<11/30$ and
\[
 \left(\frac{11}{10}\right)^{10}<\frac{13}{5},\qquad
 \left(\frac{11}{10}\right)^{23}
 <\left(\frac{13}{5}\right)^2\left(\frac{11}{10}\right)^3
 =\frac{224939}{25000}<9,
\]
where the first comparison is $25937424601<26000000000$, the desired
inequality follows. This proves \eqref{eq:coarse-strict-link} up to $\alpha=3$.

We next prove that every complementary circle arc is shorter than $b$ for
$3\le\alpha\le6$. Reflecting $A$, write its root gap in $[0,b]$ as $(L,U)$,
where
\[
 L:=-\log(1-\lambda^2),\qquad U:=-\log(1-\lambda+\lambda^2).
\]
Every other logarithmic gap has length at most
$\lambda^2(1-2\lambda)/(1-\lambda)<2/27$;
here $\lambda^2(1-2\lambda)\le1/27$ and $1-\lambda>1/2$.
Their dilations by $\alpha\le6$ therefore have length less than $4/9<b$.
The elementary bounds
\begin{equation}\label{eq:coarse-basic-bounds}
 b>\frac{91}{200},\qquad h<\frac{503}{500}
\end{equation}
will be justified explicitly below. Furthermore,
\[
 b-U=\log\left(1+\frac{\lambda^2}{1-\lambda}\right)>
 \frac{16}{93},\qquad
 U-L<\log\frac{35}{31}<\frac4{31},\qquad
 h-b<\log\frac74<\frac35.
\]
The first bound follows from $\lambda>4/11$ and
$\log(1+x)>x/(1+x)$; the last follows from
$1+3/5+(3/5)^2/2=89/50>7/4$. Hence
\[
 b(b-U)>\frac{182}{2325}>\frac{12}{155}>(U-L)(h-b),
 \qquad \frac{182}{2325}-\frac{12}{155}=\frac{62}{72075}>0.
\]
Also $\alpha b>h$ and $\alpha(U-L)<24/31<h+b$, using
$h>\log2>1/2$.
If an empty circle arc had length at least $b$, choose an empty open subarc
$I=(a,a+b)\subset(0,h)$, allowing endpoints $0,h$.
It lies in the root gap $(\alpha L,\alpha U)$, so $\alpha(U-L)\ge b$.
Its translate $I+h$ cannot partly meet $[0,\alpha b]$, since the hull endpoint
$\alpha b$ belongs to $-\alpha A$; nor can it lie wholly in that hull,
since then the same root gap would have length at least $h+b$.
Thus $a\ge\alpha b-h$, and $a+b\le\alpha U$ gives
$\alpha(b-U)\le h-b$. The last two inequalities contradict
$b(b-U)>(U-L)(h-b)$.

For $6\le\alpha\le20$, only three logarithmic gaps can be relevant.
Their lengths, in decreasing generations, are
\[
 \log\frac{1-\lambda^2}{1-\lambda+\lambda^2},\qquad
 \log\frac{(1-\lambda)(1+\lambda^2)}{1-\lambda+\lambda^3},\qquad
 \log\frac{1-\lambda^3}{1-\lambda^2+\lambda^3}.
\]
They are respectively less than
\begin{equation}\label{eq:coarse-three-gaps}
 \frac{33}{125}b,\qquad\frac{113}{1000}b,\qquad\frac{17}{200}b;
 \qquad \frac bh>\frac{453}{1000},
\end{equation}
and every deeper gap is shorter than $b/20$.
Here are explicit bounds for these claims and for
\eqref{eq:coarse-basic-bounds}. We have $c>183/500$.
The derivatives of the three logarithmic lengths have, after multiplication
by positive denominators, the numerators
\[
 1-4\lambda+\lambda^2,\quad
 \lambda(2-7\lambda+4\lambda^2-\lambda^3),\quad
 \lambda(2-6\lambda+\lambda^3).
\]
Each is negative for $183/500\le\lambda\le1/2$, since the relevant polynomial
is negative at $183/500$ and decreases thereafter. The logarithmic lengths
therefore decrease while $b$ increases.
The positive logarithmic series gives
\[
 b>2\left(\frac{183}{817}+\frac{(183/817)^3}{3}
                   +\frac{(183/817)^5}{5}\right)>\frac{45569}{100000}.
\]
For $0<z<1$,
\[
 \log\frac{1+z}{1-z}
 \le2\left(z+\frac{z^3}{3}+\frac{z^5}{5(1-z^2)}\right).
\]
Substituting $z=(R-1)/(R+1)$ gives the following rational bounds:
\[
\begin{array}{c|c}
R&2\bigl(z+z^3/3+z^5/[5(1-z^2)]\bigr)\ \text{is less than}\\ \hline
216511/191989&12021/100000\\
89866013/85378487&5123/100000\\
118871513/114383987&3849/100000
\end{array}
\]
The left column gives the three gap ratios at $\lambda=183/500$.
The entries in the right column are respectively smaller than
$33/125,113/1000,17/200$ times $45569/100000$.
Also
\[
 \sum_{k=0}^{8}\frac{(10053/10000)^k}{k!}-\frac{500}{183}
 =\frac{39624613635943016066263577308071007}
 {81984000000000000000000000000000000000}>0.
\]
Thus $h<10053/10000<503/500$, and
$(45569/100000)/(10053/10000)>453/1000$.
For the largest depth-three gap, $\log(1+x)\le x$ and
$b\ge\lambda+\lambda^2/2$ give the ratio bound
\[
 \frac{\text{gap length}}b
 \le\frac{(1-2\lambda)\lambda^2}{(1-\lambda)(1+\lambda/2)}<\frac1{20}.
\]
The rational function decreases on $[183/500,1/2]$, because its derivative
has the sign of $2-(13/2)\lambda+2\lambda^2+\lambda^3<0$.
At $183/500$ its difference from $1/20$ is
$400223/187505500>0$. The ordering by generation treats all later gaps.

Suppose now that an open circle arc of length $b$ misses $\alpha A$.
At least $\lfloor(\alpha-1)b/h\rfloor$ of its translates lie wholly in
$[-\alpha b,0]$. Each must lie in a gap of $\alpha A$.
A gap of length $\alpha g$ accommodates at most
$\lfloor(\alpha g-b)/h\rfloor+1$ such translates, and none if $\alpha g<b$.
Use the normalized hull length $t:=b/h\in(453/1000,1)$.
If $k$ of the three upper bounds in \eqref{eq:coarse-three-gaps} can be active,
the total capacity is at most
$\alpha t$ times their sum, minus $kt$, plus $k$.
The required number is greater than $(\alpha-1)t-1$.
Thus capacity is insufficient whenever
\[
 [\alpha(1-\text{sum of the active upper ratios})+k-1]t>k+1.
\]
For $6\le\alpha<1000/113$, only the first gap can be active, and
$6(1-33/125)(453/1000)-2=7/15625>0$.
For $1000/113\le\alpha<200/17$, the sum is $377/1000$.
The sufficient inequality holds if $\alpha\ge91/10$ or $t\ge47/100$;
the respective endpoint margins are
$211929/10000000$ and $173/2825$.
In the remaining rectangle, each capacity is at most one, whereas
$\lfloor(\alpha-1)t\rfloor\ge3$.
For $200/17\le\alpha\le20$, the sum is $231/500$.
The sufficient inequality holds if $\alpha\ge13$ or $t\ge49/100$;
the endpoint margins are $37141/500000$ and $173/2125$.
In the remaining rectangle the two smaller capacities are at most one
and the root capacity is at most two. If the latter is at most one,
the total is at most three, but the required number is at least four.
If it is two, then $((33/125)\alpha-1)t\ge1$.
For $\alpha\le25/2$ this gives
\[
 (\alpha-1)t\ge\frac{\alpha-1}{(33/125)\alpha-1}\ge5;
\]
for $\alpha\ge25/2$ the bound $t>453/1000$ gives the same conclusion.
The required number is therefore at least five, against capacity at most four.
All cases are impossible, proving the strict circle-gap bound.

It remains to treat $\alpha\ge20$.
For $1\le\beta<1/\lambda$, the largest complementary arc of
$\beta C_\lambda$ modulo one is at most
\[
 \begin{cases}
 \beta(1-2\lambda),&1\le\beta\le1/(1-\lambda),\\
 2-\beta,&1/(1-\lambda)\le\beta\le1+2\lambda,\\
 1-2\lambda,&1+2\lambda\le\beta<1/\lambda.
 \end{cases}
\]
Indeed, all nonroot gaps have length at most $1-2\lambda$.
A longer empty arc has a representative in $(0,1)$ lying in the root gap.
Its translate by one cannot partly meet $[0,\beta]$, because $\beta$ belongs
to the Cantor set, and cannot lie wholly in that hull, because the root gap
has length $\beta(1-2\lambda)<1$. It therefore lies outside the hull.
The first arc is consequently contained in
\[
 (\max(\beta\lambda,\beta-1),\min(\beta(1-\lambda),1)).
\]
Taking the maximum of this interval length and $1-2\lambda$ gives the
three displayed pieces. This bound decreases with $\lambda$ for fixed $\beta$.

Choose the largest integer $n\ge1$ such that
\[
 \beta:=\frac{\alpha}{h}
       \log\left(1+\frac{\lambda^n}{1-\lambda}\right)\ge1.
\]
Such an $n$ exists because the value at $n=1$ is $\alpha b/h>1$.
Concavity of the logarithm and maximality give $\beta<1/\lambda<11/4$.
The set $\alpha A$ contains the image
$\alpha\log((1-\lambda)+\lambda^n C_\lambda)$.
Compared with its secant parametrisation
$\alpha\log(1-\lambda)+\beta h C_\lambda$, every point moves in the same
direction by at most
\[
 \frac\alpha8\left(\frac{\lambda^n}{1-\lambda}\right)^2
 =\frac\alpha8(e^{\beta h/\alpha}-1)^2
 \le\frac{\beta^2(503/500)^2}
 {160(1-\beta(503/500)/20)^2}.
\]
The first estimate follows from the second derivative of $\log(1+zs)$
and $s(1-s)\le1/4$; the last uses $e^x-1\le x/(1-x)$ for $0\le x<1$.
The bound decreases when $\alpha$ increases from twenty.
Moving points in one direction by at most this amount enlarges a circle
gap by at most the same amount: remove that amount from the appropriate
end of a purported larger empty arc and pull it back.

Compare this error with $91/200$ minus $503/500$ times the preceding
three-piece gap bound at $\lambda=183/500$.
On the first piece their ratio increases with $\beta$.
On the middle piece its derivative has the sign of
\[
 2\frac{91}{200}-4\frac{503}{500}
 +\frac{503}{500}\beta+\frac{(503/500)^2}{20}\beta^2,
\]
which increases but is still
$-381191095599/312500000000<0$ at its right endpoint $\beta=433/250$.
On the last piece the ratio increases.
At the two possible maxima, the available width minus the error is
\[
 \begin{array}{c|c}
 \beta=500/317&120809903739/10800370373000>0\\
 \beta=11/4&8981604362993/74248380562500>0.
 \end{array}
\]
Hence every circle gap of $\alpha A$ is strictly shorter than $b$ for
$\alpha\ge20$. Together with the preceding ranges, this proves
\eqref{eq:coarse-strict-link} for all real $\alpha>1$.

Finally fix $\alpha=p/q>1$ in lowest terms and the integer $j$ supplied by
\eqref{eq:coarse-strict-link}. For sufficiently large integers $N$, define
on $[1-\lambda,1]$
\[
 H_N(v):=\log\left(
 \frac{1-(1-\lambda^{j+pN}v)^\alpha}
 {\alpha\lambda^{j+pN}}\right).
\]
Then $\displaystyle\lim_{N\to\infty}H_N(v)=\log v$ in $C^1$.
The thickness of $H_N((1-\lambda)+\lambda C_\lambda)$ remains greater than
one: relative to $A$, the derivative ratio of this change of coordinates
has limit one. Its hull tends to $[-b,0]$.
By \eqref{eq:coarse-strict-link}, a fixed point of $\alpha A$ remains
strictly inside its translate by $\log\alpha-jh$.
This prevents that translated Cantor set from lying in a gap of $\alpha A$.
Conversely, the diameter $\alpha b$ of $\alpha A$ exceeds the diameter of
the translated set for all large $N$, so $\alpha A$ cannot lie in one of
its bounded gaps. The strict hull linkage excludes unbounded gaps as well.
The Gap Lemma yields $u,v\in(1-\lambda)+\lambda C_\lambda$ with
\[
 \alpha\log u-H_N(v)=\log\alpha-jh.
\]
Equivalently,
\[
 x:=\lambda^{qN}u,\qquad y:=1-\lambda^{j+pN}v
 \quad\hbox{satisfy}\quad x^\alpha+y^\alpha=1.
\]
For $j+pN\ge0$ these points belong to $C_\lambda^2$ and are nontrivial.
Their $x$-coordinates lie in the pairwise disjoint intervals
$\lambda^{qN}[1-\lambda,1]$, since $\lambda^q<1-\lambda$.
Thus the intersection contains infinitely many distinct nontrivial points.
